% Pandoc template for the fed_fraud documents.
% Based on proposal_gen/templates/detailed_proposal_template/detailed_1_classic.tex
% (Classic report): centered title page, table of contents, ruled section
% headings and the Geodesic house footer. Fonts are Latin Modern (lmodern).
% Two layouts: the full report (title page and contents) and, with
% -V brief, a compact one-page-brief header in the same style.
\documentclass[letterpaper,11pt]{article}
\usepackage[T1]{fontenc}
\usepackage[utf8]{inputenc}
\usepackage{lmodern}
\usepackage[letterpaper,top=0.9in,bottom=1.2in,left=1in,right=1in,footskip=0.6in]{geometry}
\usepackage[table]{xcolor}
\usepackage{geodesicbrand}
\usepackage{amsmath,amssymb}
\usepackage{enumitem}
\usepackage{array}
\usepackage{calc}
\usepackage{longtable}
\usepackage{makecell}
\usepackage[most]{tcolorbox}
\usepackage{titlesec}
\usepackage{needspace}
\usepackage{etoolbox}
\usepackage{hyperref}
\hypersetup{colorlinks=true,urlcolor=coolblack,linkcolor=coolblack,filecolor=geoblue,
  pdftitle={Network and graph analytics},pdfauthor={\geofirm}}
\urlstyle{same}

\setlist{itemsep=2pt,leftmargin=1.4em}
\setlength{\parindent}{0pt}
\setlength{\parskip}{6pt}
\providecommand{\tightlist}{\setlength{\itemsep}{0pt}\setlength{\parskip}{0pt}}
\setcounter{secnumdepth}{0}

% Section style: house rule under each heading
\titleformat{\section}{\Large\bfseries\color{coolblack}}{\thesection}{0.6em}{}[{\color{coolblack}\titlerule[0.5pt]}]
\titleformat{\subsection}{\large\bfseries\color{geomid}}{\thesubsection}{0.6em}{}
\titleformat{\subsubsection}{\normalsize\bfseries\color{geomid}}{\thesubsubsection}{0.6em}{}
\titlespacing*{\section}{0pt}{10pt}{6pt}

% Tables are written by latex/filters.lua (navy header, banded rows)
\setlength{\LTpre}{6pt}\setlength{\LTpost}{6pt}

\begin{document}

% BRIEF HEADER
\thispagestyle{geodesic}
\begin{minipage}[c]{0.13\textwidth}\includegraphics[width=\linewidth]{icon.png}\end{minipage}\hfill
\begin{minipage}[c]{0.84\textwidth}
  {\textsc{\large \geofirm}}\par\vspace{2pt}
  {\small\color{geogray}Federal fraud expansion\enspace\textperiodcentered\enspace When: \compiletime}
\end{minipage}\par\vspace{6pt}
{\color{coolblack}\rule{\linewidth}{0.3mm}}\par\vspace{2pt}
{\LARGE\bfseries\color{coolblack} Network and graph analytics}\par
\vspace{2pt}{\large\color{geomid} Scoring rings, not single payments}\par
\vspace{-2pt}{\color{coolblack}\rule{\linewidth}{0.3mm}}\par\vspace{4pt}
\pagestyle{geodesic}

\textbf{Status:} Draft v0.1 (2026-11-04), for Steve's review. Extension
2 of 10 (see the \href{../extensions.pdf}{extensions overview}).
Reference code: \texttt{analytics/forensic/network.py}
(\texttt{communities}, \texttt{cluster\_features},
\texttt{cluster\_score}, \texttt{rank\_clusters}), tested on synthetic
data only.

\hypertarget{bottom-line}{%
\section*{Bottom line}\label{bottom-line}}
\addcontentsline{toc}{section}{Bottom line}

Organized fraud runs through groups of entities that look ordinary one
at a time: twenty hospices that each bill a plausible amount, eight PPP
borrowers with plausible payrolls, five ``small businesses'' that each
win a few set-asides. What gives them away is what they share: an
address, an officer, a phone line, a bank account, a preparer. The
forensic methods handbook already builds the shared-attribute graph and
propagates risk from known bad actors. This extension adds the step that
turns a graph into a work queue: \textbf{split the graph into
communities, describe each community with a few features, and rank
communities by a calibrated score}, so investigators open the ring, not
one member of it.

\hypertarget{graph}{%
\section*{1. Graph}\label{graph}}
\addcontentsline{toc}{section}{1. Graph}

Nodes are entities (businesses, providers, borrowers). An edge joins two
entities that share an attribute: address, officer or owner, phone,
email, bank account, registered agent, or IP or device in Phase 2. Edge
weight follows the handbook (section 8): each shared attribute adds 1 /
log(1 + degree), so a phone shared by two firms counts far more than a
registered-agent address shared by 4,000. Attributes on the benign-hub
list (registered agents, mail centers, hospitals, large payroll firms)
and attributes with more than 25 entities are dropped.

\hypertarget{from-components-to-communities}{%
\section*{2. From components to
communities}\label{from-components-to-communities}}
\addcontentsline{toc}{section}{2. From components to communities}

Connected components chain together through single weak links, so one
shared accountant can glue two unrelated rings into one blob. Weighted
label propagation splits a component into groups that are densely tied
inside and loosely tied outside. It is deterministic for a given seed,
needs no tuning and runs in near-linear time, which matters at SAM.gov
scale (several hundred thousand active registrations). Louvain or Leiden
are drop-in alternatives in a graph database.

\hypertarget{cluster-features}{%
\section*{3. Cluster features}\label{cluster-features}}
\addcontentsline{toc}{section}{3. Cluster features}

\begingroup\small\setlength{\tabcolsep}{4pt}\renewcommand{\arraystretch}{1.25}
\rowcolors{2}{white}{geolight}
\begin{longtable}[l]{@{}l>{\raggedright\arraybackslash}p{\dimexpr 0.429\linewidth-2\tabcolsep}>{\raggedright\arraybackslash}p{\dimexpr 0.429\linewidth-2\tabcolsep}@{}}
\rowcolor{coolblack}\leavevmode\color{white}\bfseries Feature & \leavevmode\color{white}\bfseries Definition & \leavevmode\color{white}\bfseries Why it matters \\
\endhead
Size & Entities in the community & Rings are usually 3 to 50 entities \\
Density & Internal edges / possible edges & Rings are tightly interconnected; franchises are star-shaped \\
Bank share & Share of members sharing a bank account with another member & The strongest single sign of common control \\
New share & Share first seen in the last two years & Rings form fast around a new program or rule change \\
Dollars & Total paid to the community & Sets priority among similar scores \\
Seed risk & Highest risk-propagation score from excluded, debarred or charged entities & Known bad actors nearby \\
\end{longtable}\endgroup

\hypertarget{score}{%
\section*{4. Score}\label{score}}
\addcontentsline{toc}{section}{4. Score}

A logistic score combines the features, with size and dollars on log
scales so a 200-outlet franchise does not outrank a six-firm ring that
shares bank accounts. The default weights in \texttt{CLUSTER\_WEIGHTS}
are starting values. Before any cluster is referred they are refit on
confirmed rings (DOJ case releases, OIG reports, the SBA case set) using
the likelihood ratio method in the handbook section 12, and the score
must meet the review plan's release bar: known rings at least five times
more common in the top 5\% than in the population.

\hypertarget{what-an-investigator-sees}{%
\section*{5. What an investigator
sees}\label{what-an-investigator-sees}}
\addcontentsline{toc}{section}{5. What an investigator sees}

For each ranked cluster the tool shows: members and their awards or
payments; the shared attributes that tie them, strongest first; a small
graph drawing; dates of formation; any seed (excluded or charged entity)
within two hops; and the reason codes behind the score. The tool preview
pages on the site follow this layout.

\hypertarget{where-it-applies-first}{%
\section*{6. Where it applies first}\label{where-it-applies-first}}
\addcontentsline{toc}{section}{6. Where it applies first}

\begingroup\small\setlength{\tabcolsep}{4pt}\renewcommand{\arraystretch}{1.25}
\rowcolors{2}{white}{geolight}
\begin{longtable}[l]{@{}>{\raggedright\arraybackslash}p{\dimexpr 0.267\linewidth-2\tabcolsep}>{\raggedright\arraybackslash}p{\dimexpr 0.444\linewidth-2\tabcolsep}>{\raggedright\arraybackslash}p{\dimexpr 0.289\linewidth-2\tabcolsep}@{}}
\rowcolor{coolblack}\leavevmode\color{white}\bfseries Roadmap idea & \leavevmode\color{white}\bfseries Graph & \leavevmode\color{white}\bfseries Seeds \\
\endhead
HHS-1, HHS-2 hospice and home health & Owners, managing employees, addresses from CMS enrollment files & LEIE, revoked providers, DOJ defendants \\
SBA-1, SBA-3 loan rings & Borrower addresses, lenders, agents, NAICS & Charged PPP and EIDL defendants \\
DOD-1, VA-1 set-aside fronts & SAM.gov officers, addresses, parent firms & SAM exclusions, SBA protests sustained \\
TREAS-2 ERC and EITC promoters & Preparers and promoters in charged cases & DOJ Tax Division cases \\
X-4 hidden common control & All of the above, linked through cross-agency matching & All exclusion lists \\
\end{longtable}\endgroup

\hypertarget{limits}{%
\section*{7. Limits}\label{limits}}
\addcontentsline{toc}{section}{7. Limits}

\begin{itemize}
\tightlist
\item
  Public data shows registrations, not control. A shared officer may be
  a part-time bookkeeper; the score asks a question, it does not answer
  one.
\item
  Rings adapt by spreading attributes thin (one phone per firm). The
  \href{red_team.pdf}{red-team playbook} sets out how the features
  respond.
\item
  Bank accounts are the strongest feature and are not public; until
  Phase 2 the bank-share feature is zero and the score leans on
  officers, phones and addresses.
\end{itemize}

\hypertarget{next-steps}{%
\section*{8. Next steps}\label{next-steps}}
\addcontentsline{toc}{section}{8. Next steps}

\begin{enumerate}
\def\labelenumi{\arabic{enumi}.}
\tightlist
\item
  Fit weights on the SBA confirmed set (weeks 5 to 6).
\item
  Run on HHS hospice enrollment data alongside HHS-1 (weeks 6 to 8).
\item
  Add a cluster view to the HHS-1 tool preview.
\end{enumerate}

\end{document}
